\section{La crisis energética y las vías de solución}

\subsection{La IA está provocando una crisis energética}

Hassabis reconoce con franqueza un hecho no del todo cómodo: \textbf{la IA está provocando una crisis energética}. Los centros de datos que requiere el entrenamiento y la inferencia a gran escala consumen una cantidad enorme de electricidad, y la demanda energética de los centros de datos a nivel mundial está creciendo a un ritmo sin precedentes.

\begin{warningbox}{La paradoja energética de la IA}
La tecnología de IA, por un lado, resuelve diversos problemas de la humanidad, pero por otro está generando un problema nuevo: una demanda enorme de energía. Si no se resuelve, el desarrollo de la IA podría desacelerarse por falta de suministro energético suficiente, o agravar el cambio climático.

Esta es una paradoja típica del tipo ``el problema que necesitas que la IA resuelva es precisamente el que la IA generó''.
\end{warningbox}

\subsection{La IA también puede resolver el problema energético}

Pero la otra cara de la moneda es: \textbf{la IA también tiene el potencial de resolver la crisis energética}. Esto incluye:

\begin{itemize}
    \item \textbf{Fusión nuclear}: la IA puede acelerar la investigación en control de plasma y ciencia de materiales, impulsando la realización de la fusión nuclear controlada
    \item \textbf{Optimización de la red eléctrica}: la gestión inteligente de la red eléctrica impulsada por IA puede mejorar de manera significativa la eficiencia energética
    \item \textbf{Descubrimiento de nuevos materiales}: la IA puede ayudar a descubrir materiales más eficientes para celdas solares, materiales de almacenamiento de energía, entre otros
    \item \textbf{Gestión energética}: la eficiencia en el uso de energía de los propios centros de datos también puede optimizarse mediante IA
\end{itemize}

\begin{importantbox}{Una perspectiva dialéctica sobre el problema energético}
La IA ha provocado un aumento acelerado en la demanda de energía, pero al mismo tiempo también podría ser la herramienta clave para resolver el problema energético. Este ciclo de ``generar el problema $\to$ resolver el problema'' es, en esencia, un rasgo común de todas las tecnologías transformadoras: la máquina de vapor creó tanto la contaminación industrial como, finalmente, impulsó el desarrollo de tecnologías de energía limpia.

Lo determinante es esto: \textbf{la velocidad de la solución debe ser mayor que la velocidad con la que se agrava el problema}. Si la IA logra impulsar un avance decisivo en la fusión nuclear dentro de 10 años, entonces el consumo energético actual será una ``inversión'' que valió la pena.
\end{importantbox}

\subsection{Resumen de la sección}

La IA está provocando una crisis energética, pero también tiene el potencial de resolver esta crisis. Lo determinante es si la velocidad de la solución puede superar la velocidad con la que se agrava el problema. La fusión nuclear, el descubrimiento de nuevos materiales y la optimización de la red eléctrica son direcciones en las que la IA puede ejercer su fuerza.

%% ==================== Section 9 ====================
